% ═══════════════════════════════════════════════════════════════
% LEHRERHINWEISE - DS-03: Repaso general
% Klasse 9, Unidad 1+2 - Klausurvorbereitung
% ═══════════════════════════════════════════════════════════════

\documentclass[11pt, a4paper]{article}

\usepackage[utf8]{inputenc}
\usepackage[T1]{fontenc}
\usepackage[ngerman]{babel}

\usepackage{helvet}
\renewcommand{\familydefault}{\sfdefault}

\usepackage[margin=2cm, top=2.5cm, bottom=2cm]{geometry}
\usepackage{enumitem}
\usepackage{tabularx}
\usepackage{booktabs}
\usepackage{array}
\usepackage{multicol}

\usepackage{fancyhdr}
\usepackage{lastpage}

\usepackage{xcolor}
\usepackage{tcolorbox}
\tcbuselibrary{skins}

% ═══════════════════════════════════════════════════════════════
% FARBEN
% ═══════════════════════════════════════════════════════════════

\definecolor{akzent}{HTML}{C41E3A}
\definecolor{hellgrau}{HTML}{F5F5F5}
\definecolor{mittelgrau}{HTML}{666666}
\definecolor{basis}{HTML}{4CAF50}
\definecolor{standard}{HTML}{2196F3}
\definecolor{erweiterung}{HTML}{9C27B0}

% ═══════════════════════════════════════════════════════════════
% BOXEN
% ═══════════════════════════════════════════════════════════════

\newtcolorbox{infobox}[1][]{
    colback=hellgrau, colframe=akzent, boxrule=1pt, arc=0pt,
    left=10pt, right=10pt, top=8pt, bottom=8pt,
    title={#1}, fonttitle=\bfseries, coltitle=white, colbacktitle=akzent
}

\newtcolorbox{diffbox}[2][]{
    colback=white, colframe=#2, boxrule=1pt, arc=0pt,
    left=8pt, right=8pt, top=6pt, bottom=6pt,
    title={#1}, fonttitle=\bfseries\small, coltitle=white, colbacktitle=#2
}

% ═══════════════════════════════════════════════════════════════
% KOPF- UND FUSSZEILE
% ═══════════════════════════════════════════════════════════════

\pagestyle{fancy}
\fancyhf{}
\renewcommand{\headrulewidth}{0.5pt}
\renewcommand{\footrulewidth}{0pt}
\lhead{\footnotesize\textsc{Lehrerhinweise} · DS-03}
\rhead{\footnotesize Repaso general · Klasse 9}
\cfoot{\footnotesize\thepage\,/\,\pageref{LastPage}}

\setlength{\parindent}{0pt}
\setlength{\parskip}{6pt}

% ═══════════════════════════════════════════════════════════════
% DOKUMENT
% ═══════════════════════════════════════════════════════════════

\begin{document}

% ═══════════════════════════════════════════════════════════════
% TITEL
% ═══════════════════════════════════════════════════════════════

\begin{center}
    {\LARGE\bfseries Lehrerhinweise: Repaso general}\\[6pt]
    {\large DS-03 · Klasse 9 · Klausurvorbereitung U1+U2}
\end{center}

\vspace{8pt}

% ═══════════════════════════════════════════════════════════════
% ÜBERBLICK
% ═══════════════════════════════════════════════════════════════

\begin{infobox}[Stundenübersicht]
\begin{tabular}{@{}ll@{}}
\textbf{Thema:} & Repaso general -- Klausurvorbereitung U1 + U2 \\
\textbf{Dauer:} & 90 Minuten (Doppelstunde) \\
\textbf{Format:} & Stationenlernen (5 Stationen) \\
\textbf{Materialien:} & AB-03, LB-03, LP-03, SLIDES-03, Lösungsblätter pro Station \\
\textbf{Voraussetzung:} & DS-01 + DS-02 abgeschlossen \\
\textbf{Klassenarbeit:} & 19./20. Januar \\
\end{tabular}
\end{infobox}

\vspace{8pt}

% ═══════════════════════════════════════════════════════════════
% FEINZIELE
% ═══════════════════════════════════════════════════════════════

\section*{Feinziele}

\begin{itemize}[leftmargin=1.5em]
    \item \textbf{FZ1:} SuS \textbf{wiederholen} alle relevanten Vokabeln aus U1 und U2 (AFB I)
    \item \textbf{FZ2:} SuS \textbf{festigen} die Bildung des pretérito perfecto (AFB I/II)
    \item \textbf{FZ3:} SuS \textbf{wenden} Objektpronomen und Bewegungsverben korrekt an (AFB II)
    \item \textbf{FZ4:} SuS \textbf{verstehen} einen klausurtypischen Text (AFB II)
    \item \textbf{FZ5:} SuS \textbf{produzieren} eigene Texte unter Klausurbedingungen (AFB III)
    \item \textbf{FZ6:} SuS \textbf{reflektieren} ihre eigene Klausurbereitschaft (Metakognition)
\end{itemize}

% ═══════════════════════════════════════════════════════════════
% STUNDENVERLAUF
% ═══════════════════════════════════════════════════════════════

\section*{Stundenverlauf}

\begin{tabularx}{\textwidth}{@{}c|l|X|c|l@{}}
\textbf{Min} & \textbf{Phase} & \textbf{Aktivität} & \textbf{SF} & \textbf{Material} \\
\midrule
5 & Einstieg & Mindmap: Klausurthemen an Tafel sammeln & PL & Tafel \\
\midrule
5 & Organisation & Stationenlernen erklären: Ablauf, Zeit, Selbstkontrolle & PL & SLIDES 2-3 \\
\midrule
\multicolumn{5}{@{}l@{}}{\textbf{Phase 2: Stationenlernen (70 Min)}} \\
\midrule
15 & Station 1 & Vokabeltraining U1+U2 (Übersetzen, Lücken) & EA & AB A1 + Lsg. \\
\midrule
14 & Station 2 & Pretérito perfecto (Konjugation, Partizipien, Sätze) & EA & AB A2 + Lsg. \\
\midrule
14 & Station 3 & Bewegungsverben + Objektpronomen (Umformung) & EA & AB A3 + Lsg. \\
\midrule
16 & Station 4 & Leseverstehen (Text + Fragen) & EA & AB A4 + Lsg. \\
\midrule
14 & Station 5 & Textproduktion (Zimmerbeschreibung) + Peer-Korrektur & EA/PA & AB A5 \\
\midrule
5 & Selbsteinsch. & Ampelkarte: Grün/Gelb/Rot je Thema & EA & AB Ampel \\
\midrule
5 & Abschluss & Offene Fragen klären, Tipps für Klausur & PL & SLIDES 9-10 \\
\bottomrule
\end{tabularx}

\textbf{Gesamtzeit:} 90 Minuten

% ═══════════════════════════════════════════════════════════════
% MATERIALIEN PRO STATION
% ═══════════════════════════════════════════════════════════════

\section*{Materialien pro Station}

\begin{tabularx}{\textwidth}{@{}c|l|X@{}}
\textbf{Nr.} & \textbf{Station} & \textbf{Material} \\
\midrule
1 & Vocabulario & AB A1 + Lösungsblatt (separat auslegen) \\
2 & Pret. perfecto & AB A2 + Regelkarte (haber-Konjugation) + Lösungsblatt \\
3 & Verbos + Pron. & AB A3 + Regelkarte (Verben/Pronomen) + Lösungsblatt \\
4 & Leseverstehen & AB A4 (enthält Text) + Lösungsblatt \\
5 & Textproduktion & AB A5 + Strukturhilfe + Peer-Korrektur-Bogen \\
\end{tabularx}

\vspace{6pt}

\textbf{Zusätzlich:} Timer (Beamer oder Uhr), Gong/Signal für Stationswechsel

% ═══════════════════════════════════════════════════════════════
% DIFFERENZIERUNG
% ═══════════════════════════════════════════════════════════════

\section*{Differenzierung}

\begin{multicols}{3}

\begin{diffbox}[{[B] Basis}]{basis}
\footnotesize
\begin{itemize}[leftmargin=1em, itemsep=2pt]
    \item Mehr Zeit an Station 1-3
    \item Regelkarten immer griffbereit
    \item Station 5: Satzanfänge nutzen
    \item Ggf. Station 4 kürzen
\end{itemize}
\end{diffbox}

\columnbreak

\begin{diffbox}[{[S] Standard}]{standard}
\footnotesize
\begin{itemize}[leftmargin=1em, itemsep=2pt]
    \item Alle Stationen vollständig
    \item Selbstständige Bearbeitung
    \item Eigenständige Fehlerkorrektur
    \item Normale Zeitvorgaben
\end{itemize}
\end{diffbox}

\columnbreak

\begin{diffbox}[{[E] Erweiterung}]{erweiterung}
\footnotesize
\begin{itemize}[leftmargin=1em, itemsep=2pt]
    \item Zusatzaufgaben bei Station 5
    \item Peer-Tutoring für [B]
    \item Komplexere Textproduktion
    \item Frühzeitig fertig: LP digital
\end{itemize}
\end{diffbox}

\end{multicols}

% ═══════════════════════════════════════════════════════════════
% LÖSUNGEN
% ═══════════════════════════════════════════════════════════════

\section*{Lösungen}

\textbf{Aufgabe 1 (Vokabeln):}
\begin{enumerate}[leftmargin=1.5em, itemsep=0pt]
    \item la nevera, 2. el dormitorio, 3. la ducha, 4. el sillón, 5. el espejo, 6. la estantería
    \item cocina, 8. salón, 9. dormitorio, 10. garaje, 11. baño, 12. cama
\end{enumerate}

\textbf{Aufgabe 2 (Pretérito perfecto):}
\begin{itemize}[leftmargin=1.5em, itemsep=0pt]
    \item A) he, has, ha, hemos, habéis, han
    \item B) hablado, comido, vivido, puesto
    \item C) Yo he hablado... / Nosotros hemos comido... / Ellos han vivido... / Tú has hecho...
\end{itemize}

\textbf{Aufgabe 3 (Pronomen):}
\begin{enumerate}[label=\alph*), leftmargin=1.5em, itemsep=0pt]
    \item La pongo en la mesa.
    \item Los subimos al piso.
    \item La sacamos del camión.
    \item Lo llevo al salón.
    \item Las bajamos del coche.
    \item La traigo de mi habitación.
\end{enumerate}

\textbf{Aufgabe 4 (Leseverstehen):}
\begin{enumerate}[label=\alph*), leftmargin=1.5em, itemsep=0pt]
    \item El piso tiene cuatro habitaciones.
    \item Le gusta porque tiene un balcón.
    \item El padre ha puesto el sofá en el salón.
    \item La madre ha puesto la ropa en el armario.
    \item Van a comprar una lámpara nueva para el salón.
\end{enumerate}

% ═══════════════════════════════════════════════════════════════
% HINWEISE
% ═══════════════════════════════════════════════════════════════

\section*{Didaktische Hinweise}

\begin{itemize}[leftmargin=1.5em]
    \item \textbf{Stationenlernen:} SuS können in beliebiger Reihenfolge arbeiten, aber empfohlen: 1→2→3→4→5
    \item \textbf{Lösungsblätter:} Separat an jeder Station auslegen (nicht ins AB integriert!)
    \item \textbf{Timer:} Sichtbar für alle, aber flexibler Wechsel (SuS arbeiten unterschiedlich schnell)
    \item \textbf{Ampelkarte:} Wichtig für Selbstreflexion -- gelbe/rote Themen vor Klausur gezielt üben
    \item \textbf{Station 5:} Peer-Korrektur motiviert und entlastet Lehrkraft
    \item \textbf{Abschluss:} Unbedingt Zeit für Fragen einplanen -- letzte Chance vor Klausur!
\end{itemize}

% ═══════════════════════════════════════════════════════════════
% AUSBLICK
% ═══════════════════════════════════════════════════════════════

\section*{Ausblick}

\textbf{Klassenarbeit:} 19./20. Januar

\textbf{Inhalte:}
\begin{itemize}[leftmargin=1.5em, itemsep=0pt]
    \item Vokabeln U1 + U2 (Alltag, Haus, Möbel)
    \item Grammatik: Pretérito perfecto, Objektpronomen, Bewegungsverben + Präpositionen
    \item Leseverstehen (ähnlich AB A4)
    \item Textproduktion (Beschreibung, Dialog, Bericht)
\end{itemize}

\textbf{Empfehlung:} SuS mit ``Rot'' bei Ampelkarte individuell beraten (Sprechstunde, zusätzliche Übungen)

\end{document}
